\documentclass[varwidth=16cm, border=2mm]{standalone}
\usepackage{amsmath}\usepackage{amssymb}\usepackage{ctex}
\usepackage{tasks}\usepackage{xcolor}\usepackage{graphicx}
\settasks{label=\Alph*., label-width=1.5em, item-indent=2em}
\begin{document}
\textbf{[0.6]} (2024·全国·高三对口高考)已知椭圆 $\frac{x^2}{9} + y^2 = 1$，过左焦点 $F$ 作倾斜角为 $\frac{\pi}{6}$ 的直线交椭圆于 $A$、$B$ 两点，则弦 $AB$ 的长为______.

\vspace{0.2cm}\par\noindent\textcolor{red}{\textbf{【答案】} 2}
\par\noindent\textcolor{blue}{\textbf{【解析】} 1. **确定椭圆的参数**\n   椭圆方程为 $\frac{x^2}{9} + y^2 = 1$。\n   根据标准方程 $\frac{x^2}{a^2} + \frac{y^2}{b^2} = 1$，我们有 $a^2 = 9$，$b^2 = 1$。\n   因此，$a = 3$，$b = 1$。\n   计算半焦距 $c$：$c^2 = a^2 - b^2 = 9 - 1 = 8$，所以 $c = \sqrt{8} = 2\sqrt{2}$。\n   椭圆的左焦点 $F$ 的坐标为 $(-c, 0) = (-2\sqrt{2}, 0)$。\n   离心率 $e = \frac{c}{a} = \frac{2\sqrt{2}}{3}$。\n\n2. **确定直线的方程**\n   直线过左焦点 $F(-2\sqrt{2}, 0)$，倾斜角为 $\frac{\pi}{6}$。\n   直线的斜率 $k = \tan(\frac{\pi}{6}) = \frac{\sqrt{3}}{3}$。\n   直线的方程为 $y - 0 = \frac{\sqrt{3}}{3} (x - (-2\sqrt{2}))$，即 $y = \frac{\sqrt{3}}{3} (x + 2\sqrt{2})$。\n\n3. **计算弦 $AB$ 的长度（方法一：焦点弦长公式）**\n   对于椭圆 $\frac{x^2}{a^2} + \frac{y^2}{b^2} = 1$，过焦点的弦长公式为 $L = \frac{2ab^2}{a^2 - c^2\cos^2\alpha}$，其中 $\alpha$ 是直线与 $x$ 轴正方向的夹角（即倾斜角）。\n   这里 $a=3, b=1, c=2\sqrt{2}, \alpha = \frac{\pi}{6}$。\n   $\cos(\frac{\pi}{6}) = \frac{\sqrt{3}}{2}$。\n   $\cos^2(\frac{\pi}{6}) = (\frac{\sqrt{3}}{2})^2 = \frac{3}{4}$。\n   代入公式：\n   $L = \frac{2 \cdot 3 \cdot 1^2}{3^2 - (2\sqrt{2})^2 \cdot \frac{3}{4}}$\n   $L = \frac{6}{9 - 8 \cdot \frac{3}{4}}$\n   $L = \frac{6}{9 - 6}$\n   $L = \frac{6}{3}$\n   $L = 2$\n\n4. **计算弦 $AB$ 的长度（方法二：韦达定理与弦长公式）**\n   将直线方程 $y = \frac{\sqrt{3}}{3} (x + 2\sqrt{2})$ 代入椭圆方程 $\frac{x^2}{9} + y^2 = 1$：\n   $\frac{x^2}{9} + (\frac{\sqrt{3}}{3} (x + 2\sqrt{2}))^2 = 1$\n   $\frac{x^2}{9} + \frac{3}{9} (x + 2\sqrt{2})^2 = 1$\n   $\frac{x^2}{9} + \frac{1}{3} (x^2 + 4\sqrt{2}x + 8) = 1$\n   两边同乘以 9，得到：\n   $x^2 + 3(x^2 + 4\sqrt{2}x + 8) = 9$\n   $x^2 + 3x^2 + 12\sqrt{2}x + 24 = 9$\n   $4x^2 + 12\sqrt{2}x + 15 = 0$\n   设 $A(x_A, y_A)$，$B(x_B, y_B)$ 是直线与椭圆的两个交点。\n   根据韦达定理：\n   $x_A + x_B = -\frac{12\sqrt{2}}{4} = -3\sqrt{2}$\n   $x_A x_B = \frac{15}{4}$\n   弦长公式为 $AB = \sqrt{(x_A - x_B)^2 + (y_A - y_B)^2}$。\n   由于直线斜率为 $k$，所以 $y_A - y_B = k(x_A - x_B)$。\n   所以 $AB = \sqrt{(x_A - x_B)^2 + k^2(x_A - x_B)^2} = \sqrt{(1+k^2)(x_A - x_B)^2} = |x_A - x_B|\sqrt{1+k^2}$。\n   计算 $(x_A - x_B)^2 = (x_A + x_B)^2 - 4x_A x_B = (-3\sqrt{2})^2 - 4(\frac{15}{4}) = 18 - 15 = 3$。\n   所以 $|x_A - x_B| = \sqrt{3}$。\n   计算 $1+k^2 = 1 + (\frac{\sqrt{3}}{3})^2 = 1 + \frac{3}{9} = 1 + \frac{1}{3} = \frac{4}{3}$。\n   所以 $\sqrt{1+k^2} = \sqrt{\frac{4}{3}} = \frac{2}{\sqrt{3}} = \frac{2\sqrt{3}}{3}$。\n   将结果代回弦长公式：\n   $AB = \sqrt{3} \cdot \frac{2\sqrt{3}}{3} = \frac{2 \cdot 3}{3} = 2$。\n\n两种方法计算结果一致，弦 $AB$ 的长为 $2$。}


\end{document}
